\documentclass[11pt]{article}
\usepackage[a4paper,margin=18mm]{geometry}
\usepackage{fontspec}
\setmainfont{Latin Modern Roman}
\usepackage{amsmath,amssymb,amsthm,amscd,graphicx,color}
\usepackage{xurl}
\usepackage[unicode,hidelinks]{hyperref}
\allowdisplaybreaks
\setlength\emergencystretch{3em}
\usepackage{enumerate}

\DeclareMathOperator{\Cliff}{Clif\/f}
\DeclareMathOperator{\Gal}{Gal}
\DeclareMathOperator{\GL}{GL}
\DeclareMathOperator{\Inn}{Inn}
\DeclareMathOperator{\Mat}{Mat}
\DeclareMathOperator{\Out}{Out}
\DeclareMathOperator{\SL}{SL}
\DeclareMathOperator{\SO}{SO}
\DeclareMathOperator{\SU}{\mathbf{SU}}
\numberwithin{equation}{section}

\newtheorem{prop}{Proposition}[section]
\newtheorem{cor}[prop]{Corollary}
\newtheorem{lem}[prop]{Lemma}
{\theoremstyle{definition}
\newtheorem{defn}[prop]{Definition}
\newtheorem{notation}[prop]{Notation}
\newtheorem{rem}[prop]{Remark}
\newtheorem{ques}[prop]{Question}
\newtheorem{eg}[prop]{Example}
\newtheorem{erratum}[prop]{Erratum}}


\makeatletter
\newcommand{\Res}[1]{\mathop{\mathbf{R}_{#1}}}

\newcommand{\ResOne}[1]{\mathop{\mathbf{R}^{(1)}_{#1}}}

\newcommand{\Lie}[1]{\mathfrak{\lowercase{#1}}}
\makeatother
\begin{document}
\begin{center}\Large Every finite-dimensional real semisimple Lie algebra has a rational form carrying rational forms of all its finite-dimensional real representations\end{center}
\noindent\textbf{Stable ID:} 2007\par
\noindent\textbf{Authors:} Dave Witte Morris\par
\noindent\textbf{Original work:} A Cohomological Proof that Real Representations of Semisimple Lie Algebras Have Q-Forms\par
\noindent\textbf{Version:} Official SIGMA 11 (2015) article 034, DOI 10.3842/SIGMA.2015.034; journal PDF and native TeX acquired 2026-10-01, exact bytes pinned by SHA256\par
\noindent\textbf{Selected task scope:} Full Proposition1.2: for every finite-dimensional real semisimple Lie algebra g\_R there exists a Lie algebra g\_Q over Q with g\_Q tensor\_Q R isomorphic to g\_R, such that every finite-dimensional real representation of g\_Q is real-conjugate to a representation over Q. This does not assert that every rational form is universal.\par
\noindent\textbf{Difficulty:} H2: The proof chooses a rational quasi-split form and modifies the twisting cocycle through central quotients so that every real-vanishing Tits obstruction also vanishes rationally. Its construction uses weak approximation for norm-one tori to lift the real cocycle, then analyzes unitary factors and central characters to establish universality simultaneously for all dominant weights.\par
\noindent\textbf{Editorial boundary note (not author text):} The full original universal rational-form Proposition 1.2 and complete Section 2 cohomological author proof are supplied. This is one simultaneous outcome for every finite-dimensional real representation of the constructed Q-form. The Tits highest-weight obstruction framework, universality criterion, norm-one-torus lifting, central-quotient cocycle construction, unitary-factor analysis and all real-vanishing/rational-vanishing obstruction comparisons are retained. Tits classification and weak approximation remain named established prerequisites. Original F/Q and torus-locus symbols are preserved with their source notice. No lattice-integrality corollary, particular-type classification or claim that every rational form is universal is selected.\par
\noindent\textbf{Source:} \url{https://sigma-journal.com/2015/034/}\quad DOI: \url{https://doi.org/10.3842/SIGMA.2015.034}\par
\noindent\textbf{License and attribution:} Copyright retained by the named authors. Published by SIGMA under CC BY 4.0: \url{https://creativecommons.org/licenses/by/4.0/}. Source: \url{https://sigma-journal.com/about.html}. Adaptation consists of standalone excerpt formatting, cover and numbering restoration; original author language and mathematical text are retained.\par
\noindent\textbf{Editorial symbol notice (not author text):} Lemma2.5 writes H1(L/F;G(L)) although the base field throughout that lemma and its proof is Q, and near the end writes T(R) where the norm-one locus is defined in T(C). The lemma opening and explicit restriction-of-scalars construction specify these objects. Original source pages remain unchanged.\par
\medskip\hrule\medskip\noindent\textbf{Author text begins below.}\par
\setcounter{section}{0}
\setcounter{subsection}{0}
\setcounter{subsubsection}{0}
\setcounter{equation}{0}
% Original source lines 77--104
\section{Introduction}

\begin{defn}
Let $\Lie G_\mathbb{Q}$ be a~Lie algebra over~$\mathbb{Q}$.
(All Lie algebras and all representations are assumed to be f\/inite-dimensional.)
\begin{enumerate}\itemsep=0pt
\item $\Lie G_\mathbb{Q}$ is \emph{universal for real representations} (or \emph{$\mathbb{R}$-universal}, for short) if
every real representation of~$\Lie G_\mathbb{Q}$ has a~$\mathbb{Q}$-form~\cite[Def\/inition~7.1]{Morris-RRepsQForms}.
This means that if $\rho \colon \Lie G_\mathbb{Q} \to \Lie{GL}(n,\mathbb{R})$ is any ($\mathbb{Q}$-linear) Lie algebra
homomorphism, then there exists $M \in \GL(n,\mathbb{R})$, such that $M \rho(x) M^{-1} \in \Lie{GL}(n,\mathbb{Q})$, for
every $x \in \Lie G_\mathbb{Q}$.
\item $\Lie G_\mathbb{Q}$ is a~\emph{$\mathbb{Q}$-form} of a~real Lie algebra~$\Lie G_\mathbb{R}$ if $\Lie G_\mathbb{Q}
\otimes_\mathbb{Q} \mathbb{R} \cong \Lie G_\mathbb{R}$.
\end{enumerate}
\end{defn}

This note uses Galois cohomology to present a~short proof of the following known result, which was f\/irst obtained~by
M.S.~Raghunathan~\cite[\S~3]{Raghunathan-ArithLattsSSGrps} in the important special case where $\Lie G_\mathbb{R}$ is
compact.

\begin{prop}[\protect{\cite[Theorem~1.2]{Morris-RRepsQForms}}]\label{HasRUniv}
Every real semisimple Lie algebra has a~$\mathbb{Q}$-form that is $\mathbb{R}$-universal.
\end{prop}

The proof in~\cite{Morris-RRepsQForms} constructs an $\mathbb{R}$-universal $\mathbb{Q}$-form explicitly, and is rather
tedious, but a~much nicer proof was given by G.~Prasad and A.~Rapinchuk~\cite[Proposition~3 and Remark~3]{PrasadRapinchuk-Prescribed}.
Assuming some fundamental results of J.~Tits~\cite{Tits-RepLinIrred}, our proof in Section~\ref{PfSect} is a~bit shorter and
more direct.
\setcounter{section}{1}
\setcounter{subsection}{0}
\setcounter{subsubsection}{0}
\setcounter{equation}{0}
% Original source lines 127--360
\section{Proof of the results stated in the introduction}\label{PfSect}

We begin by recalling a~result of J.~Tits that uses Galois cohomology to characterize the irreducible representations of
semisimple algebraic groups over f\/ields that are not algebraically closed.
However, we will state this work in the setting of Lie algebras, rather than algebraic groups.
(We deal only with semisimple Lie algebras and semisimple groups, and the f\/ields under consideration in this paper are
all of characteristic zero, so no dif\/f\/iculties arise in making this translation~\cite[Proposition~7.3.1(iii), p.~393]{DemazureGrothendieck-SGA3}.)

\begin{defn}[\protect{\cite[\S~4.2]{Tits-RepLinIrred}}] \label{BetaDefn}
Suppose $\Lie G$ is a~semisimple Lie algebra over a~subf\/ield~$F$ of~$\mathbb{C}$, and let $\mathbf{G}$ be the
corresponding simply connected, semisimple algebraic group over~$F$.
It is well known that there is a~(unique) quasi-split, simply connected algebraic group $\mathbf{G}^q$ over~$F$, and
a~$1$-cocycle $\xi \colon \Gal(\overline{F}/F) \to \overline{\mathbf{G}}{}^q$, where $\overline{\mathbf{G}}{}^q$ is the
adjoint group of~$\mathbf{G}^q$, such that $\mathbf{G}$ is~$F$-isomorphic to the Galois twist ${}^\xi \mathbf{G}^q$.
This cocycle determines a~cohomology class $[\xi] \in H^1(F; \overline{\mathbf{G}}{}^q)$.

Letting $Z(\mathbf{G}^q)$ be the center of $\mathbf{G}^q$, the short exact sequence $ e \to Z(\mathbf{G}^q) \to
\mathbf{G}^q \to \overline{\mathbf{G}}{}^q \to e $ yields a~corresponding long exact sequence of Galois cohomology sets,
including a~connecting map $\delta_* \colon H^1(F; \overline{\mathbf{G}}{}^q) \to H^2\bigl(F; Z(\mathbf{G}^q) \bigr)$.
Hence, we have a~cohomology class $\delta_* [\xi] \in H^2\bigl(F; Z(\mathbf{G}^q) \bigr)$.
(If we took a~bit more care to ensure that it is well-def\/ined, this would be the \emph{Tits class}
of~$\mathbf{G}$~\cite[p.~426]{KnusEtAl-BookOfInvols}.)

Now, f\/ix a~maximal~$F$-torus~$\mathbf{T}$ of~$\mathbf{G}^q$ that contains a~maximal~$F$-split torus, and
suppose~$\lambda$ is a~weight of~$\mathbf{T}$ that is invariant under the $*$-action of the Galois group $\Gal(\overline{F}/F)$.
Then the restriction of~$\lambda$ to $Z(\mathbf{G}^q)$ is a~$\Gal(\overline{F}/F)$-equivariant homomorphism from
$Z(\mathbf{G}^q)$ to the group~$\boldsymbol\mu$ of roots of unity in~$\mathbb{C}$, so it induces a~homomorphism
$\lambda_* \colon H^2\bigl(F; Z(\mathbf{G}^q) \bigr) \to H^2\bigl(F; \boldsymbol\mu \bigr)$.
Therefore, we may def\/ine
\begin{gather*}
\beta_{\Lie G,F}(\lambda) = \lambda_* \delta_* [\xi] \in H^2(F; \boldsymbol\mu).
\end{gather*}
\end{defn}

\begin{rem}
For a~semisimple group~$\mathbf{G}$ that is def\/ined over a~f\/ield~$F$, the def\/inition of the $*$-action of
$\Gal(\overline{F}/F)$ on the weights of a~maximal torus~$\mathbf{T}$ presupposes that $\mathbf{T}$ is def\/ined over~$F$
and contains a~maximal~$F$-split torus of~$\mathbf{G}$~\cite[p.~66]{PlatonovRapinchukBook}.
When $\mathbf{G}$ is def\/ined over~$\mathbb{Q}$, we will use the $*$-actions of both $\Gal(\overline{\mathbb{Q}}/\mathbb{Q})$ and $\Gal(\mathbb{C}/\mathbb{R})$, so we will assume that $\mathbf{T}$ is def\/ined
over~$\mathbb{Q}$ (so it is also def\/ined over~$\mathbb{R}$) and contains both a~maximal $\mathbb{Q}$-split torus and
a~maximal $\mathbb{R}$-split torus.
To see that such a~choice of~$\mathbf{T}$ is always possible, let $\mathbf{T}_1$ be any maximal $\mathbb{Q}$-split torus
of~$\mathbf{G}$.
Then there is a~maximal $\mathbb{Q}$-torus~$\mathbf{T}$ of the centralizer $\mathbf{C}_{\mathbf{G}}(\mathbf{T}_1)$ that
contains a~maximal $\mathbb{R}$-split torus~\cite[Corollary~2 of Proposition~7.8, p.~418]{PlatonovRapinchukBook}, and it is clear
that $\mathbf{T}$ has the desired properties.
\end{rem}

\begin{prop}[\protect{\cite[Theorem~7.2 and Lemma~7.4]{Tits-RepLinIrred}}]\label{HighWt}
Suppose $\Lie G$ is a~semisimple Lie algebra over a~subfield~$F$ of~$\mathbb{C}$, and~$\lambda$ is a~dominant weight.
Then:
\renewcommand{\labelenumi}{{\rm \theenumi.}}
\begin{enumerate}\itemsep=0pt
\item There is an irreducible representation ${}^F \rho_\lambda \colon \Lie G \to \Lie{GL}(n, F)$, for some~$n$, such
that \mbox{${}^F \rho_\lambda \otimes_F \mathbb{C}$} has an irreducible summand with highest weight~$\lambda$.
Furthermore, ${}^F \rho_\lambda$ is unique up to isomorphism.

\item ${}^F \rho_{\lambda_1} \cong {}^F \rho_{\lambda_2}$ if and only if~$\lambda_1$ and~$\lambda_2$ are in the same
orbit of the $*$-action of~$\Gal(\overline{F}/F)$.

\item
\label{HighWt-IrredIff}
${}^F \rho_\lambda \otimes_F \mathbb{C}$ is irreducible if and only if:
\begin{enumerate}[{\rm (a)}]\itemsep=0pt
\item
%\label{HighWt-IrredIff-Action}
$\lambda$ is invariant under the $*$-action of~$\Gal(\overline{F}/F)$, and
\item
%\label{HighWt-IrredIff-coho}
$\beta_{\Lie G,F}(\lambda)$ is the trivial element of $H^2(F; \boldsymbol\mu)$.
\end{enumerate}
\end{enumerate}
Furthermore, every~$F$-irreducible representation of~$\Lie G$ is isomorphic to ${}^F \rho_{\lambda}$, for some
dominant weight~$\lambda$.
\end{prop}

\begin{cor}
%\label{SplitOverQuadIff}
Suppose~$\Lie G$ is a~semisimple Lie algebra over~$\mathbb{Q}$, such that~$\Lie G$ splits over a~quadratic
extension of~$\mathbb{Q}$, and, for every dominant weight~$\lambda$ of~$\mathbf{T}$:
\begin{gather}
\begin{split}
& \text{if~$\lambda$ is invariant under the $*$-action of~$\Gal(\mathbb{C}/\mathbb{R})$, and $\beta_{\Lie G,\mathbb{R}}(\lambda) = 0$,}
\\
& \text{then~$\lambda$ is also invariant under the $*$-action of $\Gal(\overline{\mathbb{Q}}/\mathbb{Q})$, and
$\beta_{\Lie G,\mathbb{Q}}(\lambda) = 0$.}
\end{split}\label{SplitOverQuadIff-weight}
\end{gather}
Then $\Lie G$ is $\mathbb{R}$-universal.
\end{cor}

\begin{proof}
Since representations of $\Lie G$ are completely reducible, it suf\/f\/ices to show that every \emph{irreducible} real
representation of~$\Lie G$ has a~$\mathbb{Q}$-form.
Specif\/ically, we will show that ${}^\mathbb{Q} \rho_\lambda$ is a~$\mathbb{Q}$-form of~${}^\mathbb{R} \rho_\lambda$, for
every dominant weight~$\lambda$.

We begin by showing that if ${}^\mathbb{R} \rho_\lambda \otimes_\mathbb{R} \mathbb{C}$ is irreducible, then
${}^\mathbb{Q} \rho_\lambda \otimes_\mathbb{Q} \mathbb{C}$ is irreducible.
From Proposition~\ref{HighWt}(\ref{HighWt-IrredIff}) (with $F = \mathbb{R}$), we know that~$\lambda$ is invariant under the $*$-action
of $\Gal(\mathbb{C}/\mathbb{R})$, and that $\beta_{\Lie G,\mathbb{R}}(\lambda) = 0$.
By assumption, this implies that~$\lambda$ is invariant under the $*$-action of\/
$\Gal(\overline{\mathbb{Q}}/\mathbb{Q})$, and that $\beta_{\Lie G,\mathbb{Q}}(\lambda) = 0$.
Then, from Proposition~\ref{HighWt}(\ref{HighWt-IrredIff}) (with $F = \mathbb{Q}$), we conclude that ${}^\mathbb{Q} \rho_\lambda
\otimes_\mathbb{Q} \mathbb{C}$ is irreducible, as desired.

Since $\Lie G_\mathbb{Q}$ splits over a~quadratic extension, we know that ${}^\mathbb{Q} \rho_\lambda \otimes_\mathbb{Q}
\mathbb{C}$ is either irreducible or the direct sum of two irreducibles~\cite[Corollary~3.2(2)]{Morris-RRepsQForms}.
Therefore, ${}^\mathbb{Q} \rho_\lambda \otimes_\mathbb{Q} \mathbb{C}$ and ${}^\mathbb{R} \rho_\lambda \otimes_\mathbb{R}
\mathbb{C}$ have the same number of irreducible constituents.
(Namely, either they are both irreducible, or they are both the direct sum of $2$ irreducibles.) Since ${}^\mathbb{R}
\rho_\lambda$ is a~summand of ${}^\mathbb{Q} \rho_\lambda \otimes_\mathbb{Q} \mathbb{R}$, this implies that
${}^\mathbb{R} \rho_\lambda \cong {}^\mathbb{Q} \rho_\lambda \otimes_\mathbb{Q} \mathbb{R}$, so ${}^\mathbb{Q}
\rho_\lambda$ is a~$\mathbb{Q}$-form of ${}^\mathbb{R} \rho_\lambda$.
\end{proof}

We will also use the following variant of a~basic fact in the theory of Galois cohomology:

\begin{lem}[\protect{\cite[Theorem~5.1b, p.~77]{Kneser-LectGalCohoClassical}}]\label{KneserHasse}
If $\mathbf{G}$ is a~connected, semisimple algebraic group over~$\mathbb{Q}$, and~$L$ is any imaginary quadratic
extension of~$\mathbb{Q}$, then the natural restriction map $H^1 \bigl(L/F; \mathbf{G}(L) \bigr) \to H^1 (\mathbb{R};
\mathbf{G})$ is surjective.
\end{lem}

\begin{proof}(The author thanks A.~Rapinchuk for suggesting this argument.) Let~$\sigma$ be the nontrivial element of~$\Gal(\mathbb{C}/\mathbb{R})$.
It is well known that any cohomology class in $H^1(\mathbb{R}; \mathbf{G})$ is represented by an element~$t$ of
a~maximal torus $\mathbf{T}(\mathbb{C})$ of $\mathbf{G}(\mathbb{C})$, such that $t t^\sigma = 1$, and, since the variety
of maximal tori has weak approximation, that $\mathbf{T}$ may be chosen to be def\/ined over~$\mathbb{Q}$ (cf.\ proof
of~\cite[Proposition~6.17, p.~337]{PlatonovRapinchukBook}).

Let $\mathbf{R} = \mathrm{Res}_{L/\mathbb{Q}} \mathbf{T}$ be the torus obtained from~$\mathbf{T}$ by restriction of
scalars from~$L$ to~$\mathbb{Q}$.
Then there is a~natural isomorphism $\varphi \colon \mathbf{T}(\mathbb{C}) \to \mathbf{R}(\mathbb{R})$, such that
$\varphi \bigl(\mathbf{T}(L) \bigr) = \mathbf{R}(\mathbb{Q})$, and there is a~$\mathbb{Q}$-automorphism~$\tau$
of~$\mathbf{R}$, such that $\varphi(x^\sigma) = \varphi(x)^\tau$, for all $x \in \mathbf{T}(\mathbb{C})$.

Let $\mathbf{R}^{(1)} = \{ r \in \mathbf{R} \mid r r^\tau = 1 \}$.
We claim that this is a~subtorus of~$\mathbf{R}$.
Since $\mathbf{T}$ is def\/ined over~$\mathbb{Q}$, there is an~$L$-isomorphism $\psi \colon \mathbf{R} \to \mathbf{T}
\times \mathbf{T}$, such that $\psi \bigl(\mathbf{R}(\mathbb{R}) \bigr) = \{ (x, x^\sigma) \mid x \in
\mathbf{T}(\mathbb{C}) \}$.
Therefore, if we def\/ine $(x,y)^\pi = (y,x)$, then $\psi(r^\tau) = \psi(r)^\pi$ for all $r \in \mathbf{R}$.
So $\psi(\mathbf{R}^{(1)}) = \{ (x, x^{-1}) \mid x \in \mathbf{T} \}$ is a~subtorus of $\mathbf{T} \times \mathbf{T}$,
as claimed.

Now, $\mathbf{R}^{(1)}$ is a~torus that is def\/ined over~$\mathbb{Q}$ (since the automorphism~$\tau$ is def\/ined
over~$\mathbb{Q}$), and all $\mathbb{Q}$-tori have weak approximation at the inf\/inite place~\cite[Corollary~1 of Proposition~7.8,
p.~418]{PlatonovRapinchukBook}, so $\mathbf{R}^{(1)}(\mathbb{Q})$ is dense in $\mathbf{R}^{(1)}(\mathbb{R})$.
Hence, some $q \in \mathbf{R}^{(1)}(\mathbb{Q})$ is in the same connected component of $\mathbf{R}^{(1)}(\mathbb{R})$ as
$\varphi(t)$.
Then, letting $q' = \varphi^{-1}(q)$, we see that $q'$ is in the same connected component of $\{ w \in
\mathbf{T}(\mathbb{R}) \mid w w^\sigma = 1 \}$ as~$t$, so $q'$ represents the same cohomology class as~$t$ in
$H^1(\mathbb{R}; \mathbf{G})$.
However, since $q' \in \mathbf{T}(L)$, we see that the cohomology class of~$q'$ is in the image of $H^1 \bigl(L/F;
\mathbf{G}(L) \bigr)$, as desired.
\end{proof}

\begin{proof}[\bf Proof of Proposition~\ref{HasRUniv}] Suppose $\Lie G_\mathbb{R}$ is a~real semisimple Lie algebra, and let $\mathbf{G}$ be the
simply connected, semisimple $\mathbb{R}$-algebraic group whose Lie algebra is~$\Lie G_\mathbb{R}$.
As in Def\/inition~\ref{BetaDefn}, write $\mathbf{G} = {}^\xi \mathbf{G}^q$, where $\xi \colon \Gal(\mathbb{C}/\mathbb{R}) \to
\overline{\mathbf{G}}{}^q$ is a~$1$-cocycle and $\mathbf{G}^q$ is quasi-split.
Let $L = \mathbb{Q}[i]$.
By choosing an appropriate $\mathbb{Q}$-form, we may assume that $\mathbf{G}^q$ is a~quasi-split $\mathbb{Q}$-group that
splits over~$L$, and that the $*$-action of $\Gal(L/\mathbb{Q})$ is the same as the $*$-action of
$\Gal(\mathbb{C}/\mathbb{R})$.

Let~$\sigma$ be the nontrivial element of $\Gal(\mathbb{C}/\mathbb{R})$, f\/ix a~representative $a \in
\mathbf{G}^q(\mathbb{C})$ of $\xi(\sigma) \in \overline{\mathbf{G}}{}^q(\mathbb{C})$, and let $z = a {}^\sigma a$.
Since $\sigma^2$ is trivial and $\xi$ is a~$1$-cocycle, we know that~$z$~is trivial in $\overline{\mathbf{G}}{}^q$, so $z
\in Z(\mathbf{G}^q)(\mathbb{C})$.
This implies that~$a$~commutes with~${}^\sigma a$, so~$\sigma$ f\/ixes~$z$, which means $z \in
Z(\mathbf{G}^q)(\mathbb{R})$.

Let $\mathbf{H}$ be the product of the almost simple factors of $\mathbf{G}^q$ that are absolutely almost simple and of
type ${}^{2} A_n$ (more concretely, $\mathbf{H}$ is the product of the factors that are isomorphic to $\SU(k,\ell)$, for
some~$k$ and~$\ell$), let $Z(\mathbf{H})^2 = \{ w^2 \mid w \in Z(\mathbf{H})(\mathbb{C}) \}$, let
$\underline{\mathbf{G}}^q = \mathbf{G}^q/Z(\mathbf{H})^2$, and let $\underline z$ be the image of~$z$ in
$\underline{\mathbf{G}}^q$.
Note that $\underline{\mathbf{G}}^q$ is a~$\mathbb{Q}$-group (since $Z(\mathbf{H})^2$ is a~$\mathbb{Q}$-subgroup
of~$\mathbf{G}^q$).

We claim that we may assume $\underline z \in Z(\underline{\mathbf{G}}^q)(\mathbb{Q})$.
While proving this, we may consider each simple factor individually, so there is no harm in assuming
$\underline{\mathbf{G}}^q$ is almost simple.
This allows us to furthermore assume that $\underline{\mathbf{G}}^q$ is absolutely almost simple.
(Otherwise, since every $\mathbb{C}$-group is split, we could assume $\xi$ is trivial.) Also, since
$|{\Gal(\mathbb{C}/\mathbb{R})}| = 2$, we may assume, by replacing~$a$ with $aw$ for an appropriately chosen $w \in
\langle z \rangle$, that $|\underline z|$~is a~power of~$2$.
Assuming, as we may, that $\underline z$ is nontrivial, this implies that $\mathbf{G}^q$ is not of type ${}^{1,2} E_6$.
Then $Z(\underline{\mathbf{G}}^q)(\mathbb{R}) = Z(\underline{\mathbf{G}}^q)(\mathbb{Q})$.
(If~$\mathbf{G}^q$ is of type ${}^2 A_n$, then the def\/inition of $\underline{\mathbf{G}}^q$ implies
$|Z(\underline{\mathbf{G}}^q)| \le 2$, so every element of $Z(\underline{\mathbf{G}}^q)$ is def\/ined over~$\mathbb{Q}$.
If $\mathbf{G}^q$ is not of this type, then the desired conclusion can be verif\/ied by noting that $Z(\mathbf{G}^q)$ is
either $\boldsymbol\mu_n$, $\boldsymbol\mu_2 \times \boldsymbol\mu_2$, $\Res{L/\mathbb{Q}} \boldsymbol\mu_2$,
$\ResOne{L/\mathbb{Q}} \boldsymbol\mu_2$, or $\ResOne{L/\mathbb{Q}}
\boldsymbol\mu_4$~\cite[p.~332]{PlatonovRapinchukBook}.) This completes the proof of the claim.

The claim of the preceding paragraph implies that the cyclic subgroup $\langle \underline z \rangle$ generated
by~$\underline z$ is def\/ined over~$\mathbb{Q}$.
Hence, the quotient $\widetilde{\mathbf{G}}^q = \mathbf{G}^q/\langle z, Z(\mathbf{H})^2 \rangle$ is a~semisimple
$\mathbb{Q}$-group.
Now, since $a {}^\sigma a = z$ is trivial in $\widetilde{\mathbf{G}}^q$, we know that $\xi$ lifts to a~$1$-cocycle
$\widetilde\xi \colon \Gal(\mathbb{C}/\mathbb{R}) \to \widetilde{\mathbf{G}}^q$.
Then Lemma~\ref{KneserHasse} implies that, after replacing $\widetilde\xi$ with a~cohomologous cocycle, we may assume
$\widetilde\xi$ is the restriction of a~$1$-cocycle $\zeta \colon \Gal(L/\mathbb{Q}) \to \widetilde{\mathbf{G}}^q(L)$.
Let $\mathbf{G}_\mathbb{Q} = {}^\zeta \mathbf{G}^q$, so $\mathbf{G}_\mathbb{Q}$ is a~$\mathbb{Q}$-group that
is~$L$-split, and is isomorphic to~$\mathbf{G}$ over~$\mathbb{R}$.
Also, let $\Lie G$ be the Lie algebra of $\mathbf{G}_\mathbb{Q}$.

To complete the proof, we show that $\Lie G$ is $\mathbb{R}$-universal, by verifying~\eqref{SplitOverQuadIff-weight}.
To this end, let~$\lambda$ be a~$\Gal(\mathbb{C}/\mathbb{R})$-invariant dominant weight, such that $\beta_{\Lie
G,\mathbb{R}}(\lambda) = 0$.
Since $\mathbf{G}^q$~is~$L$-split and the $*$-action of $\Gal(L/\mathbb{Q})$ is the same as the $*$-action of
$\Gal(\mathbb{C}/\mathbb{R})$ (by the choice of the $\mathbb{Q}$-form of~$\mathbf{G}^q$), we know that~$\lambda$~is
invariant under $\Gal(\overline{\mathbb{Q}}/\mathbb{Q})$.

Since~$\zeta$ is a~$1$-cocycle into $\widetilde{\mathbf{G}}^q = \mathbf{G}^q/\langle z, Z(\mathbf{H})^2 \rangle$, we know
that, in the notation of Def\/inition~\ref{BetaDefn} with $F = \mathbb{Q}$, we have $\delta_*[\zeta] \in H^2\bigl(\mathbb{Q};
\langle z, Z(\mathbf{H})^2 \rangle\bigr)$.
Therefore, in order to show that $\beta_{\Lie G, \mathbb{Q}}(\lambda) = \lambda_* \delta_*[\zeta] = 0$, it suf\/f\/ices to
show that~$\lambda$ is trivial on both~$z$ and~$Z(\mathbf{H})^2$.
Note that, in the notation of Def\/inition~\ref{BetaDefn} with $F = \mathbb{R}$, we have $\lambda_* \delta_*[\xi] = \beta_{\Lie G,
\mathbb{R}}(\lambda) = 0$.
Under the natural identif\/ication of $H^2 \bigl(\mathbb{R}; Z(\mathbf{G}^q) \bigr)$ with
\begin{gather*}
\{ w \in Z(\mathbf{G}^q) \mid {}^\sigma w = w \} / \{ w {}^\sigma w \mid w \in Z(\mathbf{G}^q) \},
\end{gather*}
we have $\delta_* [\xi] = [z]$, so this means $\lambda(z) = 1$ (since $\omega \overline{\omega} = 1$ for all $\omega \in
\boldsymbol\mu$).
Furthermore, since the restriction of~$\lambda$ to $Z(\mathbf{G}^q)$ is a~$\Gal(\mathbb{C}/\mathbb{R})$-equivariant
homomorphism, and $Z(\mathbf{H})(\mathbb{R}) = Z(\mathbf{H})$ (cf.~\cite[p.~332]{PlatonovRapinchukBook}), we have
$\lambda \bigl(Z(\mathbf{H}) \bigr) \subseteq \boldsymbol\mu(\mathbb{R}) = \{\pm1\}$, so~$\lambda$ is also trivial on
$Z(\mathbf{H})^2$.
\end{proof}
\begin{thebibliography}{99}
\bibitem[1]{DemazureGrothendieck-SGA3}
Demazure M., Grothendieck A., Sch\'emas en groups.~III. Structure des sch\'emas
  en groupes reductifs, S\'eminaire de g\'eom\'etrie alg\'ebrique du {B}ois
  {M}arie, 1962/1964, available at
  \url{http://library.msri.org/books/sga/sga/pdf/sga3-3.pdf}.

\bibitem[3]{Kneser-LectGalCohoClassical}
Kneser M., Lectures on {G}alois cohomology of classical groups, \textit{Tata
  Institute of Fundamental Research Lectures on Mathematics}, Vol.~47, Tata
  Institute of Fundamental Research, Bombay, 1969, available at
  \url{http://www.math.tifr.res.in/~publ/ln/tifr47.pdf}.

\bibitem[4]{KnusEtAl-BookOfInvols}
Knus M.A., Merkurjev A., Rost M., Tignol J.P., The book of involutions,
  \textit{American Mathematical Society Colloquium Publications}, Vol.~44,
  Amer. Math. Soc., Providence, RI, 1998.

\bibitem[6]{Morris-RRepsQForms}
Morris D., Real representations of semisimple {L}ie algebras have {${\mathbb
  Q}$}-forms, in Algebraic Groups and Arithmetic, Tata Institute of Fundamental
  Research, Mumbai, 2004, 469--490.

\bibitem[7]{PlatonovRapinchukBook}
Platonov V., Rapinchuk A., Algebraic groups and number theory, \textit{Pure and
  Applied Mathematics}, Vol.~139, Academic Press, Inc., Boston, MA, 1994.

\bibitem[8]{PrasadRapinchuk-Prescribed}
Prasad G., Rapinchuk A., On the existence of isotropic forms of semi-simple
  algebraic groups over number f\/ields with prescribed local behavior,
  \href{http://dx.doi.org/10.1016/j.aim.2006.01.001}{\textit{Adv. Math.}} \textbf{207} (2006), 646--660.

\bibitem[9]{Raghunathan-ArithLattsSSGrps}
Raghunathan M.S., Arithmetic lattices in semisimple groups, \href{http://dx.doi.org/10.1007/BF02967980}{\textit{Proc.
  Indian Acad. Sci. Math. Sci.}} \textbf{91} (1982), 133--138.

\bibitem[11]{Tits-RepLinIrred}
Tits J., Repr\'esentations lin\'eaires irr\'eductibles d'un groupe r\'eductif
  sur un corps quelconque, \href{http://dx.doi.org/10.1515/crll.1971.247.196}{\textit{J.~Reine Angew. Math.}} \textbf{247} (1971),
  196--220.

\end{thebibliography}
\end{document}
